\scanpage{29}\mspage{28}
% L1
\uline{psTronçon} \uline{interstitiel}

% L2–3
\uline{Prop} Soit $T,T'$ \deleted{ps. tr. tels que $T<T'$}, \deleted{$a=\operatorname{Sup}T,\ b=\operatorname{Inf}T'$}, Conditions équivalentes \deleted{si $a\in{}$} \inserted{(si $a=\operatorname{Sup}T,\ b=\operatorname{Inf}T'$)}

% L4–5
a) $T\unsure{<|\prec}T'$ \inserted{et $T,T'$} ne se raccordent pas, i.e.
\[
(a\in T\text{ ou }b\in T')\implies a\neq b
\]
% L6
\deleted{b) $\exists\,x\in L$}\quad b) $\exists\,x\in L$ tel que $T\subset L_{<\unsure{x|a}},\ T'\subset L_{>\unsure{x|a}}$
% slanted
\margin{\uline{NB} b) implique déjà que $T<T'$}

% L7–9
\deleted{\unsure{b|6}}\unsure{c)|b)} $\exists$ ps. tr. $\deleted{\unsure{I}}S$ tel que
\[
T<\deleted{I}S<T'
\]
et tel que $T$ et $\deleted{I}S$, $\deleted{I}S$ et $T'$ \uline{se raccordent}.

% L10–12
De plus, ce\deleted{\unsure{t}} $\deleted{I}S$ est unique, et $\partial S=\{a,b\}$. $S$ est formé des $x\in L$ tels que
\[
T\subset L_{<a},\ T'\subset L_{>x}\qquad\text{(cf b) )}
\]

% L13–14
On appelle $S$ le ps. tr. interstitiel entre $T$ et $T'$

% L15–19
\uline{Cor} 1 Soit $T$ un tronçon, et supposons que $T$ soit minoré \inserted{\unsure{strictemt|strictement}}. Alors l'ens des \inserted{minorants} \inserted{stricts, i.e. des} $x\in L$ tels que $\unsure{\overline T|T}\subset L_{>x}$ (ens des minorants stricts) est un \deleted{tronçon} ps. tr. $S$, \deleted{que} \inserted{qui satisfait}
% L20
\begingroup
\makeatletter
\tagsleft@true
\makeatother
\[
S<T,\ S\text{ et }T\text{ se raccordent},\ S\text{ sans minorant}\tag*{\(\ast\))}
\]
\endgroup
% L21
et ces conditions \deleted{satisf} caractérisent $S$.

% L22–24
\deleted{\uline{Cor} 2 \unsure{Énoncé|Enoncé} \unsure{de|d}} On appelle $S$ le \uline{pstronçon interstitiel gauche} de \deleted{\unsure{$S$}}$T$, On définit de $=$ (si $T$ majoré) le pstr. interst. droit. de $T$.
